\documentclass{article}
\usepackage{amsmath,amssymb}
\usepackage{xurl}
\usepackage[hidelinks]{hyperref}
% Shared commands for notes in docs/paper-gaps/.

\DeclareUnicodeCharacter{2080}{\ifmmode{}_{0}\else\textsubscript{0}\fi}
\DeclareUnicodeCharacter{2081}{\ifmmode{}_{1}\else\textsubscript{1}\fi}
\DeclareUnicodeCharacter{2082}{\ifmmode{}_{2}\else\textsubscript{2}\fi}
\DeclareUnicodeCharacter{2099}{\ifmmode{}_{n}\else\textsubscript{n}\fi}

\newcommand{\Lean}{\textsc{Lean}}
\ExplSyntaxOn
\NewDocumentCommand{\leanid}{m}{
  \ifmmode
    \textsf{\detokenize{#1}}
  \else
    \group_begin:
    \urlstyle{sf}
    \tl_set:Nn \l_tmpa_tl {#1}
    \tl_replace_all:Nnn \l_tmpa_tl {\_} {_}
    \exp_args:NV \path \l_tmpa_tl
    \group_end:
  \fi
}
\ExplSyntaxOff
\providecommand{\tr}{\operatorname{tr}}
\newcommand{\tnleanIssueBase}{https://github.com/LionSR/TNLean/issues/}
\newcommand{\tnleanPrBase}{https://github.com/LionSR/TNLean/pull/}
\newcommand{\tnleanDocBase}{https://sirui-lu.com/TNLean/docs/}
\newcommand{\ghissue}[1]{\href{\tnleanIssueBase#1}{GitHub issue \##1}}
\newcommand{\ghpr}[1]{\href{\tnleanPrBase#1}{PR \##1}}
\newcommand{\doclink}[1]{\href{\tnleanDocBase#1.html}{\path{#1}}}

% Dirac notation and the identity operator, as in blueprint/src/macros/common.tex.
% \providecommand keeps note-local definitions authoritative.
\usepackage{dsfont}
\providecommand{\ket}[1]{|#1\rangle}
\providecommand{\bra}[1]{\langle #1|}
\providecommand{\braket}[2]{\langle #1 | #2 \rangle}
\providecommand{\ketbra}[2]{|#1\rangle\!\langle #2|}
\providecommand{\Id}{\mathds{1}}
\providecommand{\one}{\mathds{1}}
\providecommand{\C}{\mathbb{C}}
\providecommand{\MN}[1]{M_{#1}(\C)}
\providecommand{\MD}{M_D(\C)}

% Verdict marker: \gapnote{<kind>}{<status>}. Kind is the policy
% classification (clarification, local-correction, scope-restriction,
% unfaithful, false-source, open-gap); status is open, wip (elimination
% underway), resolved, or historical. Machine-read by the site index;
% prints a verdict line.
\newcommand{\gapnote}[2]{\par\noindent\textbf{Verdict:} #1 (#2).\par}

\title{SPT Fixed Points with Time Reversal and Reflection}
\author{}
\date{2026-10-04}
\begin{document}
\maketitle
\gapnote{scope-restriction}{resolved}

\section{What the source states}
The symmetry group in Section III.A of \cite{Cirac2021Matrix} may combine
on-site, time-reversal and reflection symmetries. Write $t,r:G\to\mathbb Z_2$
for the latter two parities and let $C_p$ be complex conjugation at odd
parity and the identity at even parity. The local symmetry is
$S_g(A)^i=\phi(g)X_g^\dagger A^iX_g$. Its phase is a one-cocycle for
$\beta=t$, whereas its virtual factor system is a two-cocycle for
$\alpha=t+r$. The source asserts that every such pair has a
zero-correlation-length representative.\footnote{Source traceability:
    \path{Papers/2011.12127/TN-Review-main.tex}, lines 1120--1130 and 1147--1157.}

\section{The resolved scope}
The earlier formalization handled linear on-site symmetries only.
\leanid{MPSTensor.mixedTensorAction} now conjugates tensor letters for time
reversal and transposes them for reflection. The ordered-product theorem
\leanid{TNLean.Algebra.symmetryLetter_list_prod} and
\leanid{MPSTensor.mpv_mixedTensorAction} show that reflection reverses the
actual finite-chain site order. Both effects are included simultaneously.

The coefficient group is the existing unitary scalar group, literally
$\mathrm U(1)$, with its conjugation action. The classes
\leanid{TNLean.Algebra.parityUnitaryH1Class} and
\leanid{TNLean.Algebra.parityUnitaryH2Class} use Mathlib's group cohomology;
they are not a second hand-built quotient or a relabeling of all complex
units. The equations are
\begin{align}
    \phi(gh) & =\phi(g)C_{t(g)}(\phi(h)),               \\
    \omega(gh,k)\omega(g,h)
             & =C_{t(g)+r(g)}(\omega(h,k))\omega(g,hk).
\end{align}
The two parities are deliberately distinct: reflection conjugates the
virtual gauge through transposition, but does not conjugate a scalar phase.

For finite $G$, the unitary regular matrices
$L_g(x,y)=\omega(g,y)\delta_{x,gy}$ realize the chosen factor system exactly.
\leanid{TNLean.Algebra.conjugateRegularUnitaryRepresentation} does not assume
that the cocycle representative is normalized. The cocycle equation itself
gives $L_gC_{t(g)+r(g)}(L_h)=\omega(g,h)L_{gh}$.

Using these virtual matrices, the dimer letters
$A^{(a,b)}=D^{-1/2}\ket a\bra b$, $D=|G|$, carry the physical action
\begin{align}
    U(g) & =\phi(g)(\overline{X_g}\otimes X_g)P^{r(g)}, \\
    \sum_j U(g)_{ij}T_{t(g),r(g)}(A^j)
         & =\phi(g)X_g^\dagger A^iX_g,
    \label{eq:rmp_spt_onsite_scope_twist}
\end{align}
where $P$ swaps the two dimer factors and $T_{t,r}$ is entrywise conjugation
followed by transpose when the corresponding parity is odd.
\leanid{MPSTensor.exists_mixed_spt_fixedPoint} proves injectivity, the
rank-one idempotent transfer map, physical unitarity, the semilinear group
law and the prescribed tensor symmetry. The physical matrices satisfy
$U(gh)=U(g)C_{t(g)}(U(h))$; conjugation is applied separately when the
symmetry is antiunitary. The finite-chain vector has phase $\phi(g)^N$, as
proved by \leanid{MPSTensor.mixedSptAction_mpv}.

\leanid{MPSTensor.parityUnitaryH2Class_eq_of_mixedSptSymmetry} proves that
any other unitary virtual realization of the same physical symmetry and
character has the same $\mathrm U(1)$ cohomology class. Injectivity gives
pointwise scalar gauge equivalence, unitarity forces unit-modulus scalars,
and their ratio is the twisted coboundary.

\section{Conventions and remaining boundaries}
At $t=r=0$ the local action and cocycle equations specialize to the on-site
ones. The existing arbitrary-character on-site theorems remain available.
The construction retains the corrected dimer tensor: it does not reinstate
the erroneous printed tensor discussed in
\cite{gap:rmp_spt_fixed_point_tensor}. Its bond space is nonzero because
$D=|G|$ for a finite group. The result is the fixed-point realization and its
class independence; it does not by itself prove a new mixed-symmetry gapped
interpolation theorem.

\bibliographystyle{alpha}
\bibliography{references}
\end{document}
